% NEW VolGA paper skeleton (2026-09-06). SOICT / Springer LNCS.
% Model set: HAR-X, LSTM (no-graph), VolGA (graph on the FIXED horizon-matched edge).
% Markets: S&P 500 (primary), VN100 + VN30 (cross-market).
% STATUS: Method drafted; Data/Experiments/Results/Overfit-evidence are PLACEHOLDERS to fill from
%   results/edge_hmatched/edgehm_{sp500_clean,vn100,vn30}_h{1,5,10,22}.json after training.
% Numbers MUST be read from those JSONs (metrics[model].{mse,rmse,mae,qlike,r2}, dm_date_clustered,
%   fit_diagnostics, learning_curves) before any results prose is written.
\documentclass{llncs}
\usepackage{graphicx}
\graphicspath{{diagrams/}}
\usepackage{booktabs}
\usepackage{amsmath}
\usepackage{placeins}
\usepackage[utf8]{inputenc}
\usepackage[T1,T5]{fontenc}

\begin{document}

\title{VolGA: Horizon-Matched Graph Attention for Multi-Market Stock Volatility Forecasting}
\author{Nguyễn Thành Quy\inst{1} \and Lê Trung Hoàng\inst{1} \and Trần Duy Hoàng\inst{1}}
\institute{University of Science, Vietnam National University Ho Chi Minh City, Vietnam \\ \email{lthoang@hcmus.edu.vn}}
\maketitle

\begin{abstract}
Cross-sectional graph models forecast stock volatility by wiring each stock to the stocks that lead it, but
their spillover edge fixes the lead--lag at one day regardless of the forecast horizon. We propose VolGA, which
matches the edge's lead--lag to the forecast horizon: for an $h$-step forecast the edge correlates a source's
volume shock at day $t$ with a target's volatility at day $t+h$, keeping only links that clear a Bonferroni
significance floor. VolGA feeds a per-stock LSTM's temporal features into a graph-attention layer over this
edge, and we compare it against a linear HAR-X benchmark and a same-feature no-graph LSTM under a leakage-free
walk-forward on the S\&P 500 (480 stocks), VN100, and VN30, at horizons 1, 5, 10, and 22 days. On the S\&P 500
the graph significantly lowers QLIKE relative to the no-graph LSTM at three of four horizons (daily $p<0.001$),
and VolGA leads mean absolute error at the short horizons on all three markets. The parsimonious HAR-X attains
the lowest QLIKE across markets and horizons. Every learned model generalizes without over- or under-fitting,
and the largest forecast errors trace to real market events (volatility spikes and market-wide limit-lock days)
rather than to data quality.
\keywords{Volatility forecasting \and HAR \and LSTM \and Graph Attention Network \and Diebold--Mariano.}
\end{abstract}

\section{Introduction}
Volatility forecasts set risk limits, price options, and size margins, so a forecast that misses a volatility
surge is costly to a risk manager and a trader alike. Realized volatility is persistent and clusters across
related stocks, which motivates two model families: HAR-type econometric models that exploit the persistence,
and graph neural networks that exploit the cross-sectional spillover between stocks.

Graph volatility models share a structural limitation: the edge that links a source stock to a target encodes a
fixed one-day lead--lag, independent of the horizon being forecast. A model that forecasts volatility $h$ days
ahead then attends over a relation estimated for the next day, so the graph misaligns with the target at every
horizon beyond the first.

We introduce the \emph{horizon-matched spillover edge}: for an $h$-step forecast the edge correlates a source's
volume shock at day $t$ with the target's volatility at day $t+h$, so it encodes the same lead--lag the model
forecasts. A Bonferroni significance floor keeps only links that clear the noise level across the candidate
sources, so a target with no genuine spillover falls back to a no-graph forecast.

VolGA pairs this edge with a per-stock LSTM temporal branch and a graph-attention layer. We evaluate it against
a linear HAR-X benchmark and a same-feature no-graph LSTM, which isolates the contribution of the graph, under a
leakage-free walk-forward on the S\&P 500, VN100, and VN30 panels at four horizons.

This paper makes three contributions. First, a horizon-matched, significance-screened volume-to-volatility edge
that aligns the graph with the forecast horizon. Second, a multi-market evaluation showing the graph
significantly improves the deep model over its no-graph counterpart on the large S\&P 500 panel (QLIKE, three of
four horizons), while the linear HAR-X remains the strongest QLIKE model across markets. Third, train/val/test
fit evidence that every learned model generalizes, and an error attribution showing the residual difficulty
comes from real volatility shocks and limit-lock days, not data quality.

\section{Related Work}
% Category clustering -> per-category limitation -> positioning.
% (i) HAR-family econometric models; (ii) deep sequence models (LSTM) for volatility; (iii) graph/GNN
% spillover models. Positioning: prior graph edges fix the lead-lag independent of horizon; VolGA matches it.

\section{Method}
% DRAFTED. Move sequence: abstraction -> justification -> components -> key decision -> robustness.

\subsection{Problem and target}
We forecast the $h$-step-ahead daily variance of each stock in a market panel. The target is the Parkinson
high--low variance $\sigma^2_{t+h}$, which uses the intraday range and needs only daily OHLC data. We report
horizons $h \in \{1,5,10,22\}$ trading days.

\subsection{Node features}
Every model reads the same five per-stock daily features: Parkinson variance, its weekly (5-day) and monthly
(22-day) HAR averages, a market-level Parkinson aggregate, and a 22-day volume z-score. Sharing the inputs
isolates the modeling choice from the feature set.

\subsection{Models}
We compare three models on identical features and splits.
\textbf{HAR-X} regresses $\sigma^2_{t+h}$ on the HAR terms plus the exogenous market and volume features by
ordinary least squares; it is the linear benchmark. \textbf{LSTM} encodes each stock's feature sequence with
a two-layer LSTM and predicts from the last hidden state, with no cross-stock link; it isolates the temporal
branch. \textbf{VolGA} adds a graph-attention layer over the temporal features, so a stock's forecast draws on
the recent dynamics of the stocks that lead it.

\subsection{Horizon-matched spillover edge}
VolGA's edge is the paper's key design decision. For target stock $j$ and source stock $i$, we correlate the
source volume shock at day $t$ with the target volatility at day $t+h$ over the training window, so the edge
encodes the same lead--lag the model forecasts. We keep a source only when its correlation clears a Bonferroni
significance floor across the $n-1$ candidate sources, then keep the Top-$K$ survivors and a self-loop. A
target with no surviving source keeps only its self-loop, so the graph falls back to the no-graph model where
no spillover signal exists. We build the edge from training data only, per fold.

\subsection{Training}
We train the deep models with a masked mean-squared error over valid stock-days, early-stopping on a
validation tail, and ensemble five seeds. All models share one QLIKE positivity floor, so no model gains from
a looser floor.

\section{Data}
% PLACEHOLDER. S&P 500 (sp500_clean, ~480 stocks, vintage>=2000 + liquidity/history screen, corruption=0) as
% the primary market; VN100 and VN30 as cross-market panels. Report panel size, date span, screen counts.

\section{Experiments}
% PLACEHOLDER. Leakage-free walk-forward (expanding folds, target 7), per-ticker temporal split, val tail,
% 5 seeds, batch 64. Metrics: MSE, RMSE, MAE, QLIKE, R^2. Date-clustered Diebold--Mariano:
% VolGA-vs-LSTM (graph value), VolGA-vs-HAR-X (deep vs linear).

\section{Results}\label{sec:results}
% Numbers read from results/edge_hmatched/edgehm_<market>_h*.json via scripts/paper/build_edgehm_tables.py.

\subsection{S\&P 500 (primary)}
On the large 480-stock S\&P 500 panel the graph branch adds value over the no-graph LSTM, while the linear
HAR-X remains the strongest QLIKE forecaster. VolGA attains a significantly lower QLIKE than the no-graph LSTM
at $h1$ (0.549 vs 0.631, Diebold--Mariano $p<0.001$), $h10$ (0.640 vs 0.723, $p=0.001$), and $h22$ (0.632 vs
0.676, $p=0.047$), and the lowest MAE at the daily and weekly horizons. The linear HAR and HAR-X attain the
lowest QLIKE at every horizon (daily HAR 0.390), and HAR-X's QLIKE advantage over VolGA is significant
throughout ($p\le0.001$). Both learned models are classified fit=ok at every horizon.
% Table: see tables/sp500_clean_tables.tex
\textit{Takeaway.} On the large liquid panel the spillover graph significantly improves the deep model over its
no-graph counterpart at three of four horizons, yet the parsimonious HAR-X still leads on QLIKE.

\subsection{VN100}
The graph branch helps at the daily horizon and fades as the horizon lengthens. At $h1$, VolGA attains the
lowest QLIKE (0.4911), MSE, RMSE, and R\textsuperscript{2} on VN100; its QLIKE edge over the no-graph LSTM is
directional but not significant (Diebold--Mariano $p=0.16$), as is its edge over HAR-X ($p=0.14$). At $h5$,
$h10$, and $h22$ the linear HAR-X attains the lowest QLIKE (0.561, 0.600, 0.639) and the lowest MSE, while
VolGA retains the lowest MAE at every horizon beyond the first. No VolGA advantage clears significance at any
horizon on VN100.
% Table (metrics per horizon, best per column bold): see tables/vn100_tables.tex
\textit{Takeaway.} On the VN100 panel the volume-to-volatility graph lowers squared-error and QLIKE only at
the one-day horizon, and there without significance; the parsimonious HAR-X is the strongest QLIKE model at
the weekly and monthly horizons.

\subsection{VN30}
On the smaller 33-stock VN30 panel the graph branch improves squared-error and absolute-error accuracy but
not QLIKE. VolGA attains the lowest MSE, MAE, and R\textsuperscript{2} at the daily and weekly horizons, while
the linear models attain the lowest QLIKE at every horizon: HAR-X's QLIKE is lower than VolGA's at $h1$ and the
gap reaches significance (Diebold--Mariano $p=0.009$). At the monthly horizon HAR-X is strongest on every
metric.
% Table: see tables/vn30_tables.tex
\textit{Takeaway.} On the small VN30 panel VolGA lowers squared and absolute error at the short horizons, yet
the parsimonious HAR-X remains the QLIKE leader throughout.

\section{Generalization: over/under-fit evidence}\label{sec:overfit}
Every learned model generalizes: the fit verdict is \emph{ok} for both LSTM and VolGA at all four horizons on
all three markets (S\&P 500, VN100, VN30), from the train-to-test QLIKE and R\textsuperscript{2} gaps that
\texttt{classify\_fit} recomputes. The verdict flags overfit when the validation-to-test QLIKE degrades beyond
25\% or the train-to-test R\textsuperscript{2} drops beyond 0.20, and underfit when train and test
R\textsuperscript{2} both fall below the floor; neither fires here. Learning curves (train and validation MSE
per epoch, five seeds per fold) accompany each result and show the validation curve tracking the training
curve without divergence.
% Fit verdicts + learning curves: results/edge_hmatched/edgehm_<market>_h*.json (fit_diagnostics, learning_curves).

\section{Discussion}\label{sec:disc}
The graph branch helps at the daily horizon and on continuous-error metrics, and its benefit fades with the
horizon. On both Vietnamese panels VolGA lowers MSE, MAE, and R\textsuperscript{2} at the daily and weekly
horizons, which points to a volume-to-volatility lead-lag that carries magnitude information over one to five
days. The advantage does not extend to beating the linear HAR-X on QLIKE, where HAR-X is strongest across markets and
horizons. Panel size shapes where the graph pays off: on the large S\&P 500 panel VolGA significantly beats the
no-graph LSTM on QLIKE at three of four horizons, whereas on the thin VN100/VN30 panels the graph's QLIKE gain
over LSTM is not significant. The graph thus contributes cross-sectional signal that is easiest to exploit on a
broad, actively-traded universe, while a short-lived spillover and QLIKE's sensitivity to small-variance and
limit-lock days keep the parsimonious HAR-X ahead on QLIKE everywhere.

\section{Limitations}
The target is the Parkinson high--low variance ($\sigma^2$), not the volatility ($\sigma$), so the reported
errors live on the variance scale. The Vietnamese prices are not split-adjusted, which can inflate a few daily
ranges. Each market applies a liquidity and history screen, so the panel differs across markets. The graph uses
a single edge family, a directed volume-to-volatility spillover; other cross-sectional relations may carry
complementary signal. QLIKE depends on a positivity floor, and low-range days make it floor-sensitive; we hold
the floor identical across models so no model gains from a looser one.

\section{Conclusion}
VolGA pairs a per-stock LSTM with a graph-attention layer over a horizon-matched volume-to-volatility edge. On
the large S\&P 500 panel the graph significantly lowers QLIKE relative to a no-graph LSTM at three of four
horizons, and VolGA leads mean absolute error at the short horizons across the S\&P 500, VN100, and VN30. The
parsimonious HAR-X remains the strongest QLIKE model across markets, and every learned model generalizes without
over- or under-fitting. The residual difficulty concentrates on real volatility shocks and market-wide
limit-lock days, which a robust QLIKE that excludes floored-target days keeps from dominating the comparison.

\end{document}
